\section{Research question and hypotheses}
\label{sec:question}

\Cref{sec:problem} established the premise: a domain reconstructed from public text is thin in
predictable ways. This section turns it into a testable question, three hypotheses and a falsifier.

\begin{keyquestionbox}
  \textbf{RQ-B.} Can AI reconstruct what is written about a domain, and predict where its own
  reconstruction is thin, so that the location of the human knowledge residual is identified before
  any human is asked?
\end{keyquestionbox}

RQ-B is the reoriented project's central question (\repo{REORIENTATION.md}, \S1, \S22). It asks
whether the \emph{same} system that reconstructs a domain can say where the reconstruction is
weakest, in a way that lines up with where a competent human would add something. \evint{} A
\enquote{no} would not make reconstruction worthless; it would mean targeting adds nothing over
asking experts everything, or over simpler signals such as text density.

\subsection{Hypotheses}

\paragraph{H1 --- Evidence reconstruction with provenance is feasible and checkable.} A system can
extract atomic claims from public sources and tie every evidence item to the exact span it came
from. It labels each claim's evidence status; claims without a locatable span stay labeled
\code{synthetic_extrapolation} and never act as a criterion. These links are checked rather than
trusted. H1 includes an adversarial part: gating the reconstruction lets fewer planted
falsehoods and fewer false answers through than the ungated model.

\paragraph{H2 --- Methodology lenses plus an absence check predict the human residual better than
the strongest baseline.} Constructs from cognitive task analysis and related methods
(\cref{sec:foundations}) are applied as lenses over the reconstructed evidence. Each lens firing is
followed by a check of whether the element the lens expects is stated there (the absence check,
\emph{closure} in the code). Together they rank areas by where an expert would add something the
evidence lacks. They do so better than the strongest baseline without this machinery, such as text
density or a plain language-model guess.

\paragraph{H3 --- Prioritized questions reduce expert time.} Questions chosen from the gap map
recover at least 1.25 times as many validated, previously unwritten items per expert hour as an
untargeted interview of equal length. The 1.25 is the smallest effect of interest (SESOI), fixed
before data, and a null result counts against H3. The PoC's breadth ranking and selection by expected information gain (not built)
are tested as separate selectors (\cref{sec:validation}, V6).

\evint{} H2 is the central hypothesis; \repo{REORIENTATION.md} and the program notes call it
\enquote{H1}. The order is one of dependency: H2 means something only if H1 holds, and H3 only if
H2's ranking beats chance. A negative result at one link does not falsify the link below it.

\subsection{The falsifier}

\begin{evfutbox}
  H2 is falsified if the upper bound of a 90\% confidence interval on $\Delta$AUROC falls below
  a SESOI $\delta$, fixed before any gold is acquired. $\Delta$AUROC is the gap map's area-level
  AUROC for ranking residual-present areas above others, minus that of the strongest
  pre-registered baseline. The
  blind prospective study (V5) reads this rule and decides H2. It has not run.
\end{evfutbox}

The gap map must beat its strongest competitor by at least $\delta$, not just chance. The test can
show that this margin is absent only when enough areas exist: about 150 per class for
$\delta = 0.05$. A single CTA gold gives perhaps 15--25 areas and cannot. \Cref{sec:validation}
therefore splits the work. A retrospective study pooled over several published CTA golds
\emph{screens} the gap map and releases money for the prospective study only if the map looks
promising. The prospective study is sized for about 150 areas per class and \emph{decides} H2.
Before either, a closure study tests whether the absence check works at all; it needs no grant.

\evhyp{} One further hypothesis sits outside this chain. \textbf{H-OSS:} the gap map, adapted to
code and commit evidence, predicts where knowledge is lost when developers leave. It concerns
organizational artifacts, not the human residual, so its result can neither support nor refute
H2.

No hypothesis has yet been tested against a human criterion (\cref{tab:hypothesis-status}).

\begin{table}[htbp]
  \centering\small
  \caption{What the zero-participant PoC could test for each hypothesis, and its status after N2 and
    N3. No cell reports a result against human knowledge.}
  \label{tab:hypothesis-status}
  \begin{tabularx}{\linewidth}{@{}p{0.05\linewidth}>{\raggedright\arraybackslash}p{0.33\linewidth}
      >{\raggedright\arraybackslash}X@{}}
    \toprule
    Hyp. & What the PoC could test & Status after the PoC \\
    \midrule
    H1 & Whether a pipeline with a claim ledger, verbatim provenance and cross-model verification
      runs end to end on real sources. &
      \evobs\ Engineering part: yes (\cref{sec:n2}). Adversarial part: N2's registered verdict is
      INCONCLUSIVE; E-PLANT's Continue branch closed on the validity floor, and the gated arms were
      never scored. \\
    H2 & Whether lenses fire on real text and whether the absence check separates \enquote{unstated}
      from \enquote{stated elsewhere}. &
      \evobs\ Mechanism only. Lenses fire on three ledgers; the judge's verdict does not depend on
      which evidence it is shown, so the absence check is not yet informative (\cref{sec:n3}). No
      $\Delta$AUROC has been computed. Next test: the closure study (V1), before any grant. \\
    H3 & Whether ranked questions can be produced at all. &
      \evfut\ Not tested. Template questions exist; no expert has answered one. \\
    \bottomrule
  \end{tabularx}
\end{table}
